% ==============================================================================
% PAQUETES Y CONFIGURACIÓN GENERAL (APA 7)
% ==============================================================================

\usepackage[utf8]{inputenc}
\usepackage[spanish, es-tabla]{babel} % es-tabla asegura "Tabla" en lugar de "Cuadro"
\usepackage{csquotes}
\usepackage{graphicx}
\usepackage{booktabs}  % Tablas elegantes según estándar APA (sin líneas verticales)
\usepackage{amsmath}
\usepackage{ragged2e} % Permite forzar justificación plena (\justifying)
\usepackage{indentfirst} % Sangría de primera línea en TODO párrafo, incluido el primero tras cada encabezado
\usepackage{hyperref}

% Configuración de citas y bibliografía APA 7.ª Edición
\usepackage[style=apa, backend=biber]{biblatex}
\addbibresource{referencias.bib}

% Enlaces configurados en color negro para un aspecto tradicional impreso.
% Los enlaces siguen siendo funcionales en el PDF, pero sin tinta de color
% (puntos de guía en el índice, sin resaltado azul).
\hypersetup{
    colorlinks=true,
    linkcolor=black,
    citecolor=black,
    urlcolor=black
}

% Líneas punteadas (puntos de guía) en el Índice General para los títulos de
% nivel 1. El article.cls actual define \l@section sin puntos; se restaura el
% formato clásico con \@dottedtocline conservando la alineación a la izquierda.
\makeatletter
\renewcommand*\l@section{\@dottedtocline{1}{0em}{2.3em}}
\makeatother
